% Daniel B. Dimitrov — CV, tailored to GSK Associate Director, Translational Data Sciences (445321)
% Compile with: xelatex DanielDimitrov_CV.tex  (run twice)
\documentclass[11pt,a4paper]{article}

\usepackage[top=1.15cm,bottom=1.15cm,left=1.5cm,right=1.5cm]{geometry}
\usepackage[default]{sourcesanspro}
\usepackage[dvipsnames]{xcolor}
\usepackage{graphicx}
\usepackage{enumitem}
\usepackage{titlesec}
\usepackage{fontawesome5}
\usepackage{tabularx}
\usepackage[hidelinks]{hyperref}
\usepackage{paracol}
\usepackage{setspace}
\setstretch{1.0}

% ---------- palette ----------
\definecolor{accent}{HTML}{0B6E70}   % deep teal
\definecolor{inktext}{HTML}{2B2B2B}
\definecolor{softgrey}{HTML}{6A6A6A}

\color{inktext}
\setlength{\parindent}{0pt}
\pagestyle{empty}
\hyphenpenalty=600
\emergencystretch=1.5em
\hyphenation{bio-log-i-cal om-ics multi-omics}

% ---------- section styling ----------
\titleformat{\section}
  {\color{accent}\large\bfseries\scshape}
  {}{0em}{}[{\color{accent}\vspace{-2pt}\titlerule[0.6pt]}]
\titlespacing*{\section}{0pt}{9pt}{5pt}

% ---------- helpers ----------
\newcommand{\cvrole}[3]{% role, org+location, dates
  {\bfseries #1}\hfill{\color{softgrey}\itshape #3}\\
  {\color{softgrey}#2}\par\vspace{2pt}}
\newcommand{\cvline}[3]{% one-line role
  \par\noindent{\bfseries #1}, {\color{softgrey}#2}\hfill{\color{softgrey}\itshape #3}\par\vspace{4pt}}
\newcommand{\sep}{\;{\color{accent}\textbullet}\;}
\newenvironment{cvitems}
  {\begin{itemize}[leftmargin=1.05em,itemsep=2pt,parsep=0pt,topsep=2.5pt,label={\color{accent}\small\textbullet}]}
  {\end{itemize}\vspace{5pt}}
\newcommand{\pub}[1]{\par\hangindent=1.1em\hangafter=1 #1\vspace{6pt}}
\newcommand{\me}{\textbf{Dimitrov,~D.}}
\newcommand{\pkg}[2]{\href{#2}{\color{accent}#1}}

\begin{document}

% ================= HEADER =================
\begin{minipage}[c]{\textwidth}
{\Huge\bfseries\color{accent}DANIEL B.\ DIMITROV}\\[4pt]
{\large\color{softgrey}Computational Biologist --- Multi-omics, Spatial \& Clinical Data for Target and Biomarker Discovery}\\[7pt]
{\small
\faEnvelope\; \href{mailto:daniel.boykov.dimitrov@gmail.com}{daniel.boykov.dimitrov@gmail.com}\sep
\faPhone*\; +49\,177\,705\,6912\sep
\faMapMarker*\; Heidelberg, Germany\\[3pt]
\faGithub\; \href{https://github.com/dbdimitrov}{github.com/dbdimitrov}\sep
\faGraduationCap\; \href{https://tinyurl.com/scholarddimitrov}{Google Scholar}
}
\end{minipage}

\vspace{16pt}
I build computational frameworks that turn single-cell, spatial and multi-omics data --- integrated
with clinical and phenotypic data --- into quantitative evidence on disease mechanism, biomarkers and
targets. Shipped as production-quality open-source software: 300,000+ installs, used across academia and industry.
\vspace{4pt}

\columnratio{0.66}
\setlength{\columnsep}{20pt}
\begin{paracol}{2}

% ================= MAIN COLUMN =================
\section{Experience}

\cvrole{Postdoctoral Researcher}{Stegle Lab, EMBL \& DKFZ, Heidelberg, DE}{Sep 2024 -- Present}
\begin{cvitems}
  \item Developing probabilistic models to infer perturbed and niche-specific
        intercellular signalling in spatial omics data --- surfacing signalling axes
        as candidate targets
  \item Conceived and direct 5 parallel projects, incl.\ \textbf{EHRx} --- interpreting transformer
        foundation models of \textbf{electronic health records}: steering internal features to expose
        encoded comorbidity patterns, tested for \textbf{polygenic-risk enrichment against UK~Biobank
        genetic and clinical data}; mentored students to (co-)first authorship
\end{cvitems}

\cvrole{Early Stage Researcher (PhD)}{Saez Lab, UKHD, Heidelberg, DE}{Sep 2020 -- Aug 2024}
\begin{cvitems}
  \item Developed methods to model cell-cell communication from single-cell and
        spatial (multi-)omics data; established current best practices in the field
  \item Built and maintain LIANA \& LIANA+ (\textit{Nature Cell Biology}) ---
        300,000+ installs --- with continuous integration, unit testing, and
        clear documentation
  \item Applied these to \textbf{patient cohorts} with clinical collaborators --- CKD and IgA
        nephropathy (MSCA \textit{STRATEGY-CKD}), and the Human Cell Atlas \textit{Nature} study of
        inflammatory gut disease, where my analysis pinpointed inflammatory-niche signalling axes
\end{cvitems}

\cvrole{Visiting Predoc}{Zeller Group, EMBL --- multi-modal host-microbiome data}{Sep 2022 -- Mar 2023}
\vspace{2pt}
\cvrole{Visitor}{Glorieux Lab, Ghent --- transcriptomics \& proteomics, CKD patient cohorts}{May -- Jun 2022}
\vspace{2pt}
\cvrole{Junior Data Engineer}{Experian, Sofia --- ETL pipeline for terabytes of data}{Dec 2019 -- Aug 2020}
\vspace{2pt}
{\color{softgrey}Earlier: MSc thesis, Glasgow --- published R app for RNA-Seq analysis (2019)\sep IMPRS internship, MPI-CBG Dresden --- cryo-EM (2018)}

\section{Education}

\cvrole{PhD in Computational Biology}{Heidelberg University}{2020 -- 2024}
\vspace{2pt}
\cvrole{MSc in Bioinformatics --- Distinction}{University of Glasgow}{2018 -- 2019}
\vspace{2pt}
\cvrole{BSc (Hons) in Cell Biology --- First Class}{University of Stirling}{2014 -- 2018}

% ================= SIDEBAR =================
\switchcolumn
\raggedright

\section{Programming}
\textbf{Python} and \textbf{R}\\[3pt]
\textbf{ML:} PyTorch, PyG, scikit-learn\\[3pt]
\textbf{Omics:} scverse, BioConductor\\[3pt]
\textbf{Engineering:} git, Continuous Integration, Unit testing, conda, docker, Agile development\\[3pt]
\textbf{Main developer:}\\
\pkg{LIANA+}{https://github.com/saezlab/liana-py},
\pkg{LIANA}{https://github.com/saezlab/liana},
\pkg{Cellina}{https://github.com/PMBio/cellina}\\[3pt]
\textbf{Co-developer:}\\
\pkg{decoupler}{https://github.com/saezlab/decoupler-py},
\pkg{MetalinksDB}{https://metalinks.omnipathdb.org/}

\section{Consortia}
scverse\\[3pt]
Single-cell Best Practices\\[3pt]
Open Problems in Single-cell~Analysis

\section{Talks \& Events}
\textbf{Invited:}\\
VIB Inflammation Research, Ghent 2026\\[3pt]
GAF Summer School, Sofia 2025\\[3pt]
MOPITAS Symposium, SDU 2024\\[5pt]
\textbf{Selected \& institutional:}\\
AI \& Biology @ EMBL 2026\\[3pt]
AI Seminar @ EMBL 2026\\[3pt]
VIB Spatial Omics 2024\\[5pt]
\textbf{Co-organised:}\\
scverse Hackathon 2023\\[3pt]
ECCB 2022 tutorial

\section{Funding \& Awards}
MSCA ITN Fellowship (STRATEGY-CKD)\\[3pt]
CZI grant --- named key person\\[3pt]
Research-Based Learning Prize

\section{Languages}
English (fluent), German (intermediate), Bulgarian (native)

\end{paracol}

% ================= PAGE 2 =================
\newpage

\section{Selected Publications and Preprints \normalsize\mdseries(out of 18 \textbullet\ 4,000+ citations, h-index 14 \textbullet\ \# co-first, * corresponding)}

\pub{Moeed, A., Schrod, S., Rohbeck, M., Bonder, M.J., Lutsik, P., Stegle, O.* and \me*, 2026.
Querying Counterfactuals on Tissue Graphs with Supervised Disentanglement (Cellina).
\textit{arXiv preprint}.}

\pub{\textbf{Dimitrov,~D.}\#{}*, Schrod, S.\#{}*, Rohbeck, M. and Stegle, O.*, 2025.
Interpretation, extrapolation and perturbation of single cells.
\textit{Nature Reviews Genetics}.}

\pub{Farr, E.B.\#, \me\#, Turei, D., Schmidt, C., Lobentanzer, S., Dugourd, A. and Saez-Rodriguez, J., 2024.
MetalinksDB: a flexible and contextualizable resource of metabolite-protein interactions.
\textit{Briefings in Bioinformatics}.}

\pub{Oliver, A.J., Huang, N., et al.\ incl.\ \me, 2024.
Single-cell integration reveals metaplasia in inflammatory gut diseases.
\textit{Nature} (Human Cell Atlas).}

\pub{Sch\"afer, P.S.L., \me, Villablanca, E.J. and Saez-Rodriguez, J., 2024.
Integrating single-cell multi-omics and prior biological knowledge for a functional
characterization of the immune system. \textit{Nature Immunology}.}

\pub{\me, Sch\"afer, P.S.L., Farr, E., Rodriguez Mier, P., Lobentanzer, S., Dugourd, A., et al.\
and Saez-Rodriguez, J., 2024. LIANA+ provides an all-in-one framework for cell--cell
communication inference. \textit{Nature Cell Biology}.}

\pub{Baghdassarian, H.\#, \me\#, Armingol, E.\#, Saez-Rodriguez, J. and Lewis, N.E., 2024.
Combining LIANA and Tensor-cell2cell to decipher cell-cell communication across
multiple samples. \textit{Cell Reports Methods}.}

\pub{\me, T\"urei, D., Garrido-Rodriguez, M., Burmedi, P.L., Nagai, J.S., Boys, C.,
Ramirez Flores, R.O., Kim, H., Szalai, B., Costa, I.G., Valdeolivas, A. and Saez-Rodriguez, J., 2022.
Comparison of methods and resources for cell-cell communication inference from
single-cell RNA-Seq data. \textit{Nature Communications}, 13(1).}

\pub{B\"ulow, R.D.\#, \me\#, Boor, P. and Saez-Rodriguez, J., 2021.
How will artificial intelligence and bioinformatics change our understanding of
IgA Nephropathy in the next decade? \textit{Seminars in Immunopathology}, 43(5).}

\section{Supervision \& Leadership}
\begin{cvitems}
  \item 10 students supervised to date (5 PhD, 3 MSc, 2 BSc), 5 currently active. Four ongoing projects
        I conceived and direct: \textbf{EHRx} (EHR foundation-model interpretability; UK~Biobank
        polygenic-risk enrichment), \textbf{LIANA++} (multimodal cell--cell communication),
        \textbf{KIARA} (niche-specific perturbation-effect modelling), and \textbf{Cellina$\to$imaging}
        (pathology foundation models for tissue counterfactuals)
  \item Mentored students through to authorship --- first author on Cellina (under my senior and
        co-corresponding authorship) and co-first on MetalinksDB --- and routinely direct student
        work into released, tested codebases
  \item Cross-functional and community leadership: Human Cell Atlas, Single-cell Best Practices,
        and Open Problems in Single-Cell Analysis (cell--cell communication task);
        co-organised the scverse Hackathon 2023 and the ECCB 2022 tutorial
\end{cvitems}

\section{References}

\begin{tabularx}{\textwidth}{@{}XX@{}}
\textbf{Dr.\ Julio Saez-Rodriguez} & \textbf{Dr.\ Oliver Stegle}\\
PhD Supervisor & Postdoctoral Supervisor\\
Head of Research, EMBL-EBI & Acting Head of AI, EMBL \& DKFZ\\
\href{mailto:julio.saez@uni-heidelberg.de}{julio.saez@uni-heidelberg.de} &
\href{mailto:oliver.stegle@embl.de}{oliver.stegle@embl.de}\\
\end{tabularx}

\end{document}
